% ============================================================
% ECON 308 — International Trade
% Lecture 1 Quiz — Student Version
% Fall 2026  |  Muhebullah Karimzada  |  UNBC
% ============================================================
\documentclass[11pt]{article}
\usepackage[margin=1in,top=1.2in,bottom=1.2in]{geometry}
\usepackage[utf8]{inputenc}
\usepackage[T1]{fontenc}
\usepackage{lmodern}
\usepackage{amsmath,amssymb,bm}
\usepackage{booktabs,array,multirow,adjustbox}
\usepackage{enumitem}
\usepackage{xcolor}
\usepackage{fancyhdr}
%\setlength{\headheight}{22pt}
\addtolength{\topmargin}{-10pt}
\usepackage{tcolorbox}
\usepackage{titlesec}
\usepackage{hyperref}
\tcbuselibrary{skins}
\hypersetup{colorlinks=false,hidelinks}

%% ── Colours ────────────────────────────────────────────────
\definecolor{UNBCDark}  {RGB}{0,0,0}
\definecolor{UNBCGold}  {RGB}{0,0,0}
\definecolor{SoftGray}  {RGB}{244,246,248}

%% ── Header / Footer ────────────────────────────────────────
\pagestyle{fancy}
\fancyhf{}
\renewcommand{\headrulewidth}{0pt}

\lhead{\small\textbf{\color{UNBCDark}Name:\enspace\underline{\hspace{6cm}}}}
\rhead{\small\color{UNBCDark}Student ID:\enspace\underline{\hspace{4.5cm}}}
\cfoot{\small\thepage}

%% ── Section style ──────────────────────────────────────────
\titleformat{\section}[block]
  {\large\bfseries\color{UNBCDark}}
  {}
  {0em}
  {}
  [{\color{UNBCGold}\titlerule[0.8pt]}]
\titlespacing*{\section}{0pt}{1.2em}{0.6em}

%% ── Custom boxes ───────────────────────────────────────────
\newtcolorbox{infobox}{
  colback=SoftGray, colframe=UNBCDark,
  left=5mm, right=5mm, top=4mm, bottom=4mm,
  boxrule=0.7pt, arc=3mm}

\newtcolorbox{instructbox}{
  colback=UNBCDark!5, colframe=UNBCDark!50,
  left=4mm, right=4mm, top=3mm, bottom=3mm,
  boxrule=0.5pt, arc=2mm}

\newtcolorbox{databox}{
  colback=SoftGray, colframe=UNBCDark!40,
  left=4mm, right=4mm, top=3mm, bottom=3mm,
  boxrule=0.4pt, arc=2mm}

%% ── Answer space macro ─────────────────────────────────────
\newcommand{\answerspace}[1]{\vspace{#1}\noindent\color{UNBCDark!20}\hrule\vspace{0.3em}}

% ============================================================
\begin{document}
% ============================================================

%% ── Title block ─────────────────────────────────────────────
\begin{center}
{\LARGE\bfseries\color{UNBCDark}ECON 308 \;—\; International Economic Relations}\\[0.5em]
{\large\bfseries Quiz 1\quad\textcolor{UNBCGold}{$|$}\quad September 21, 2026}\\[0.2em]
{\small University of Northern British Columbia\quad
 Instructor: M.\,Karimzada}
\end{center}

\vspace{0.4em}

\vspace{0.6em}



\vspace{0.6em}



%% ════════════════════════════════════════════════════════════
\section{Part A:\enspace Multiple Choice\quad\normalfont\small
         (5 questions $\times$ 2 points = 10 points)}
%% ════════════════════════════════════════════════════════════



\bigskip

%% ── Question 1 ───────────────────────────────────────────────

\begin{enumerate}


\item In the expenditure approach to GDP, Canada's national income identity is
\[Y \;=\; C + I + G + (X - M).\]
Suppose exports are \$800 billion and imports are \$840 billion, with all other
components (C, I, G) held constant. Which statement \emph{best} describes
Canada's trade balance and its \emph{direct} effect on GDP?

\medskip
\begin{enumerate}[label=\textbf{\Alph*)},leftmargin=2em,itemsep=0.4em]
  \item Trade \emph{surplus} of \$40 billion; this adds \$40 billion to GDP.
  \item Trade \emph{deficit} of \$40 billion; this reduces GDP by \$40 billion.
  \item Trade \emph{deficit} of \$40 billion; this reduces GDP by \$80 billion
        because both exports and imports enter the identity separately.
  \item Trade \emph{surplus} of \$840 billion; imports create income for trading
        partners and therefore do not reduce domestic GDP.
\end{enumerate}



\bigskip\bigskip

\item The simple gravity model predicts bilateral trade as
\[T_{ij} \;=\; A\cdot\frac{Y_i\cdot Y_j}{D_{ij}},\]
where $\alpha = \beta = \gamma = 1$ and $A = 1$. Suppose Country $i$'s GDP
increases by 20\%, while Country $j$'s GDP and the bilateral distance $D_{ij}$
remain unchanged. Predicted bilateral trade $T_{ij}$ will:

\medskip
\begin{enumerate}[label=\textbf{\Alph*)},leftmargin=2em,itemsep=0.4em]
  \item Increase by exactly 20\%.
  \item Increase by 40\%, since GDP growth compounds across both the exporting
        and importing directions.
  \item Increase by approximately 44\%, due to a multiplier effect from economic
        expansion.
  \item Decrease by 20\%, as larger economies become more self-sufficient and
        require fewer imports.
\end{enumerate}



\bigskip\bigskip

%% ── Question 3 ───────────────────────────────────────────────
\item McCallum (1995) compared interprovincial trade within Canada to
Canada--U.S.\ trade at the same physical distance, \emph{after} the Canada--U.S.\
Free Trade Agreement (CUSFTA) had already eliminated most bilateral tariffs. He
found that Canadian provinces traded approximately \textbf{20 times more} with
each other than with equally distant U.S.\ states. The most accurate economic
conclusion from this finding is:

\medskip
\begin{enumerate}[label=\textbf{\Alph*)},leftmargin=2em,itemsep=0.4em]
  \item Physical transport costs are the primary remaining barrier to Canada--U.S.\ trade.
  \item Canadian interprovincial trade data systematically overstates actual trade
        volumes due to measurement differences.
  \item National borders impose substantial non-tariff costs on trade, even under
        free trade agreements.
  \item The CUSFTA agreement failed to liberalize Canada--U.S.\ merchandise trade
        in any meaningful way.
\end{enumerate}

\medskip

%% ── Question 4 ───────────────────────────────────────────────
\item According to Lecture 1, which of the following \emph{best} describes the
historical pattern of world trade as a share of global output since 1870?

\medskip
\begin{enumerate}[label=\textbf{\Alph*)},leftmargin=2em,itemsep=0.4em]
  \item World trade as a share of global output rose continuously from 1870 to
        the present, driven by steady improvements in transportation technology.
  \item World trade expanded from 1870 to 1914, then collapsed between 1914 and
        1945 due to war and protectionism, before recovering and expanding again
        after World War II.
  \item World trade was negligible before 1945 and only became significant after
        the invention of containerization in the late 1950s.
  \item World trade grew rapidly from 1870 to 1914, then recovered immediately
        after World War I and continued rising without interruption through to
        the present.
\end{enumerate}


\medskip

\item The ``Dutch Disease'' mechanism, as discussed in Lecture 1, describes
how a resource-sector boom can harm the manufacturing sector. Which of the
following sequences \emph{best} describes this mechanism in the Canadian context
(2002--2014)?

\medskip
\begin{enumerate}[label=\textbf{\Alph*)},leftmargin=2em,itemsep=0.4em]
  \item A commodity boom diverts capital investment away from manufacturing,
        directly reducing manufacturing output through factor-market reallocation.
  \item A commodity boom causes the Bank of Canada to raise interest rates,
        increasing financing costs and depressing manufacturing investment.
  \item A commodity boom improves the trade surplus so dramatically that
        foreign-produced manufactures flood the domestic market, eliminating
        domestic manufacturing.
  \item A commodity boom raises foreign demand for the Canadian dollar, causing
        the CAD to appreciate, which makes Canadian manufactured exports more
        expensive for foreign buyers and reduces manufacturing competitiveness.
\end{enumerate}
\end{enumerate}

\bigskip



%% ════════════════════════════════════════════════════════════
\section{Part B:\enspace Numerical Problem\quad\normalfont\small
         (10 points)}
%% ════════════════════════════════════════════════════════════


\bigskip

\noindent\textbf{The Gravity Model and Canada's Trade Policy}

\medskip
\noindent Canada's GDP is \textbf{\$2,400 billion}. Consider Canada's trade
relationships with three partners:

\vspace{0.5em}
\begin{center}
\begin{tabular}{lccc}
\toprule
\textbf{Partner} & \textbf{GDP (\$B)} & \textbf{Distance from Canada (km)}\\
\midrule
Australia      & 1{,}600  & 12{,}000 \\
Mexico         & 1{,}200  & \phantom{1}3{,}000 \\
United States  & 25{,}000 & \phantom{1}1{,}000 \\
\bottomrule
\end{tabular}
\end{center}
\vspace{0.5em}

\noindent Use the following gravity model \emph{throughout} this problem:
\[T_{ij} \;=\; A\cdot\frac{Y_i\cdot Y_j}{D_{ij}},\qquad
  A = 1\quad\text{(values in \$billions and km)}\]

\bigskip

%% (a)
\noindent\textbf{(a)}\quad(2 points)\enspace
Compute the gravity model's predicted bilateral trade for each country pair:
$T_{\text{CA}}$ (Canada--Australia), $T_{\text{CM}}$ (Canada--Mexico), and
$T_{\text{CU}}$ (Canada--United States).

\vspace{.5cm}

%% (b)
\noindent\textbf{(b)}\quad(2 points)\enspace
Calculate the ratio $T_{\text{CU}}\,/\,T_{\text{CA}}$.
Decompose this ratio into the two gravity-model factors---\emph{size} and
\emph{distance}---that explain why Canada trades so much more with the U.S.\
than with Australia. Interpret each factor.

\vspace{.5cm}

%% (c)
\noindent\textbf{(c)}\quad(2 points)\enspace
CPTPP membership reduces Canada's effective distance to Australia by
\textbf{25\%}, from 12{,}000\,km to 9{,}000\,km. Compute the new predicted
trade $T_{\text{CA}}^{\text{CPTPP}}$ and the percentage increase relative to
the baseline in part~(a).

\vspace{.5cm}

%% (d)
\noindent\textbf{(d)}\quad(4 points)\enspace
Suppose Canada's initial GDP is $Y = \$2{,}400$ billion, with net exports
$\text{NX} = +\$48$ billion. A U.S.\ recession causes Canadian exports to fall
by \textbf{\$96 billion}, while imports ($M$), consumption ($C$), investment
($I$), and government spending ($G$) all remain unchanged. Using the expenditure
identity $Y = C + I + G + \text{NX}$:
\begin{enumerate}[label=(\roman*),leftmargin=2em,itemsep=1.5em]
  \item Calculate the new level of net exports after the shock.
  \item Calculate the new level of Canadian GDP.
\end{enumerate}






\vspace{1em}
{\color{UNBCGold}\hrule height 0.6pt}
\begin{center}
\small\itshape
--- End of Quiz ---\qquad Total: 20 points\qquad Good luck!
\end{center}

\end{document}
